%File: anonymous-submission-latex-2024.tex
\documentclass[letterpaper]{article} % DO NOT CHANGE THIS
\usepackage[submission]{aaai24}  % DO NOT CHANGE THIS
\usepackage{amsmath}
\usepackage{amsthm}
\usepackage{times}  % DO NOT CHANGE THIS
\usepackage{helvet}  % DO NOT CHANGE THIS
\usepackage{courier}  % DO NOT CHANGE THIS
\usepackage[hyphens]{url}  % DO NOT CHANGE THIS
\usepackage{graphicx} % DO NOT CHANGE THIS
\urlstyle{rm} % DO NOT CHANGE THIS
\def\UrlFont{\rm}  % DO NOT CHANGE THIS
\usepackage{natbib}  % DO NOT CHANGE THIS AND DO NOT ADD ANY OPTIONS TO IT
\usepackage{caption} % DO NOT CHANGE THIS AND DO NOT ADD ANY OPTIONS TO IT
\frenchspacing  % DO NOT CHANGE THIS
\setlength{\pdfpagewidth}{8.5in} % DO NOT CHANGE THIS
\setlength{\pdfpageheight}{11in} % DO NOT CHANGE THIS
%
% These are recommended to typeset algorithms but not required. See the subsubsection on algorithms. Remove them if you don't have algorithms in your paper.
\usepackage{algorithm}
\usepackage{algorithmic}
\theoremstyle{definition}
\newtheorem{definition}{Definition}
\newtheorem{proposition}{Proposition}
\usepackage{subcaption}
\usepackage{amssymb}

%
% These are are recommended to typeset listings but not required. See the subsubsection on listing. Remove this block if you don't have listings in your paper.
\usepackage{newfloat}
\usepackage{listings}
\DeclareCaptionStyle{ruled}{labelfont=normalfont,labelsep=colon,strut=off} % DO NOT CHANGE THIS
\lstset{%
	basicstyle={\footnotesize\ttfamily},% footnotesize acceptable for monospace
	numbers=left,numberstyle=\footnotesize,xleftmargin=2em,% show line numbers, remove this entire line if you don't want the numbers.
	aboveskip=0pt,belowskip=0pt,%
	showstringspaces=false,tabsize=2,breaklines=true}
\floatstyle{ruled}
\newfloat{listing}{tb}{lst}{}
\floatname{listing}{Listing}
%
% Keep the \pdfinfo as shown here. There's no need
% for you to add the /Title and /Author tags.
\pdfinfo{
/TemplateVersion (2024.1)
}

% DISALLOWED PACKAGES
% \usepackage{authblk} -- This package is specifically forbidden
% \usepackage{balance} -- This package is specifically forbidden
% \usepackage{color (if used in text)
% \usepackage{CJK} -- This package is specifically forbidden
% \usepackage{float} -- This package is specifically forbidden
% \usepackage{flushend} -- This package is specifically forbidden
% \usepackage{fontenc} -- This package is specifically forbidden
% \usepackage{fullpage} -- This package is specifically forbidden
% \usepackage{geometry} -- This package is specifically forbidden
% \usepackage{grffile} -- This package is specifically forbidden
% \usepackage{hyperref} -- This package is specifically forbidden
% \usepackage{navigator} -- This package is specifically forbidden
% (or any other package that embeds links such as navigator or hyperref)
% \indentfirst} -- This package is specifically forbidden
% \layout} -- This package is specifically forbidden
% \multicol} -- This package is specifically forbidden
% \nameref} -- This package is specifically forbidden
% \usepackage{savetrees} -- This package is specifically forbidden
% \usepackage{setspace} -- This package is specifically forbidden
% \usepackage{stfloats} -- This package is specifically forbidden
% \usepackage{tabu} -- This package is specifically forbidden
% \usepackage{titlesec} -- This package is specifically forbidden
% \usepackage{tocbibind} -- This package is specifically forbidden
% \usepackage{ulem} -- This package is specifically forbidden
% \usepackage{wrapfig} -- This package is specifically forbidden
% DISALLOWED COMMANDS
% \nocopyright -- Your paper will not be published if you use this command
% \addtolength -- This command may not be used
% \balance -- This command may not be used
% \baselinestretch -- Your paper will not be published if you use this command
% \clearpage -- No page breaks of any kind may be used for the final version of your paper
% \columnsep -- This command may not be used
% \newpage -- No page breaks of any kind may be used for the final version of your paper
% \pagebreak -- No page breaks of any kind may be used for the final version of your paperr
% \pagestyle -- This command may not be used
% \tiny -- This is not an acceptable font size.
% \vspace{- -- No negative value may be used in proximity of a caption, figure, table, section, subsection, subsubsection, or reference
% \vskip{- -- No negative value may be used to alter spacing above or below a caption, figure, table, section, subsection, subsubsection, or reference

\setcounter{secnumdepth}{1} %May be changed to 1 or 2 if section numbers are desired.

% The file aaai24.sty is the style file for AAAI Press
% proceedings, working notes, and technical reports.
%

% Title

% Your title must be in mixed case, not sentence case.
% That means all verbs (including short verbs like be, is, using,and go),
% nouns, adverbs, adjectives should be capitalized, including both words in hyphenated terms, while
% articles, conjunctions, and prepositions are lower case unless they
% directly follow a colon or long dash
\title{Supplementary materials for Leveraging Imperfect Symmetry for Multi-Agent Reinforcement Learning}
\author{
    %Authors
    % All authors must be in the same font size and format.
    Written by AAAI Press Staff\textsuperscript{\rm 1}\thanks{With help from the AAAI Publications Committee.}\\
    AAAI Style Contributions by Pater Patel Schneider,
    Sunil Issar,\\
    J. Scott Penberthy,
    George Ferguson,
    Hans Guesgen,
    Francisco Cruz\equalcontrib,
    Marc Pujol-Gonzalez\equalcontrib
}
\affiliations{
    %Afiliations
    \textsuperscript{\rm 1}Association for the Advancement of Artificial Intelligence\\
    % If you have multiple authors and multiple affiliations
    % use superscripts in text and roman font to identify them.
    % For example,

    % Sunil Issar\textsuperscript{\rm 2},
    % J. Scott Penberthy\textsuperscript{\rm 3},
    % George Ferguson\textsuperscript{\rm 4},
    % Hans Guesgen\textsuperscript{\rm 5}
    % Note that the comma should be placed after the superscript

    1900 Embarcadero Road, Suite 101\\
    Palo Alto, California 94303-3310 USA\\
    % email address must be in roman text type, not monospace or sans serif
    proceedings-questions@aaai.org
%
% See more examples next
}

%Example, Single Author, ->> remove \iffalse,\fi and place them surrounding AAAI title to use it
\iffalse
\title{My Publication Title --- Single Author}
\author {
    Author Name
}
\affiliations{
    Affiliation\\
    Affiliation Line 2\\
    name@example.com
}
\fi

\iffalse
%Example, Multiple Authors, ->> remove \iffalse,\fi and place them surrounding AAAI title to use it
\title{My Publication Title --- Multiple Authors}
\author {
    % Authors
    First Author Name\textsuperscript{\rm 1},
    Second Author Name\textsuperscript{\rm 2},
    Third Author Name\textsuperscript{\rm 1}
}
\affiliations {
    % Affiliations
    \textsuperscript{\rm 1}Affiliation 1\\
    \textsuperscript{\rm 2}Affiliation 2\\
    firstAuthor@affiliation1.com, secondAuthor@affilation2.com, thirdAuthor@affiliation1.com
}
\fi


% REMOVE THIS: bibentry
% This is only needed to show inline citations in the guidelines document. You should not need it and can safely delete it.
\usepackage{bibentry}
% END REMOVE bibentry

\begin{document}

\maketitle


\section{Proof of the Proposition 1}

We first recap that the Partially Symmetric Markov game $\mathcal{M}_g$, which has the following two defined properties:

%\begin{definition}[Partially Symmetric Markov game]
%The partial Symmetry Markov game $\mathcal{M}_g=(N, S,  \left\{A_{i}\right\}_{i=1}^{n}, R, T, \Psi, g, \epsilon,  \delta )$ is a cooperative Markov game that satisfies the conditions of partial reward invariance and partial transition invariance. The state and action transformation are defined as $L_{g}: S \rightarrow S$ and $K_{g}: A \rightarrow A$, respectively. For state-action pairs $(s,a) \in \Psi$, we denote the transformed state-action pairs as $(gs,ga)$. Here, $gs=L_{g}(s)$ and $ga=K_{g}(a)$ for short. 
The partial reward invariance is characterized with
\begin{equation}
    |R(s, a) - R(gs, ga)| \leq \epsilon \label{rewardequ}.
\end{equation}
The partial transition invariance is characterized using the Maximum Mean Discrepancy (MMD) between the distributions of transitions $T(s'|s, a)$ and $T(gs'|gs, ga)$: 
\begin{equation}
    \text{MMD}_{\mathcal{F}}\left(T(s'|s, a), T(gs'|gs, ga)\right)\leq \delta \label{transequ}.
\end{equation}
Please note that $T(gs'|gs, ga)$ is also a distribution of $s'$ with the sampling process $s' \sim T(gs'|gs, ga)$ being defined as 1)~$gs' \sim T(gs'|gs, ga)$ and then 2) $s'=L_g^{-1}(gs')$.

The MMD is a measure of distance used to quantify the discrepancy between two distributions under a general class of  bounded mapping functions $f \in \mathcal{F}$. Eq.(2) can be expressed as:
\begin{equation*}
\begin{aligned}
    &\text{MMD}_{\mathcal{F}}\left(T(s'|s, a), T(gs'|gs, ga)\right) \\
    &= \sup_{f \in \mathcal{F}} \left| \mathbb{E}_{s' \sim T(s'|s, a)}[f(s')] - \mathbb{E}_{s' \sim T(gs'|gs, ga)}[f(s')] \right| \\   
\end{aligned}.
\end{equation*}
%\end{definition}





%After recapping the Partially Symmetric Markov game, we provide the proof of proposition 1 in the following. We will use mathematical induction to prove this proposition.
\begin{proposition}[Performance Error Bound]
\label{prop1}
\textit{If a Partial Symmetry Markov game $\mathcal{M}_g$  satisfies the conditions in equations \eqref{rewardequ} and \eqref{transequ} for all $(s, a) \in \Psi$, then the performance error $\textit{Error}_{\mathcal{M}_g}=|Q^{\star}(s, a) - Q^{\star}(gs, ga)|$ is bounded by $\frac{\epsilon}{1-\gamma} + \frac{\gamma \delta}{1-\gamma}$.}
\end{proposition} \label{prop1}


\begin{proof} We will apply the mathematical induction method for the proof. We define the $m$-step optimal discounted action value function recursively for all $(s, a) \in \Psi$ and for all non-negative integers $m$ as follows:
\begin{equation}  
\begin{aligned} 
    &Q_{m}(s, a)=R(s, a)+\gamma \sum_{s^{\prime} \in S}[T(s'|s, a) \max _{a^{\prime} \in A} Q_{m-1}(s^{\prime}, a^{\prime})], \label{eq1}
\end{aligned}
\end{equation}
where we set $ Q_{-1}\left(s^{\prime}, a^{\prime}\right)=0 $. Please note that because the reward $R$ is bounded, $Q_{m}(s, a)$ is bounded for any natural number $m$, and so do the functions of $\max_{a'} Q_m(s', a')$ and $\max_{a'} Q_m(gs', ga')$. Therefore, they are all elements of $\mathcal{F}$.

%We can rewrite the equation $Q_m(s,a)$ as follows:
%\begin{equation*} 
%    Q_{m}(s, a)=R(s, a)+\gamma \sum_{s^{\prime} \in S}[T(s'|s, a) V_{m-1}(s^{\prime})].
%\end{equation*}

%Because the reward is bounded, there exists a constant $\Delta$ such that $R(s,a)+\Delta>0$. Therefore, we can assume the reward $R(s,a)$ to be positive and conduct the proof without the loss of generality. This implies that $Q_m$ are $V_m$ are all positive by default.
\noindent
\textbf{Step 1 (base step)}: For the base case of $m=0$, we have
{\scriptsize
\begin{equation*}
|Q_{0}(s, a) - Q_{0}(gs, ga)|  = |R(s,a) - R(gs,ga)|  \leq \epsilon
\end{equation*}}

\textbf{Step 2 (base step)}: For the base case of $m=1$, we have
{\scriptsize
\begin{equation*}
\begin{aligned} 
& \left|Q_1(s, a) - Q_1(gs, ga)\right| \\
& = \Biggl |R(s, a) + \gamma \sum_{s^{\prime}} T(s'|s, a)  \max_{a^{\prime}} Q_0(s^{\prime}, a^{\prime}) \\
& \quad - R(gs, ga) - \gamma \sum_{s^{\prime}} T(gs'|gs, ga)  \max_{a^{\prime}} Q_0(gs^{\prime}, ga^{\prime}) \Biggl | \\
& \leq \Biggl |R(s, a)- R(gs, ga) \Biggl | + \gamma \Biggl |\sum_{s^{\prime}} T(s'|s, a)  \max_{a^{\prime}} Q_0(s^{\prime}, a^{\prime}) \\
& \quad  -  \sum_{s^{\prime}} T(gs'|gs, ga)  \max_{a^{\prime}} Q_0(gs^{\prime}, ga^{\prime}) \Biggl | \\
& \leq \epsilon + \gamma \Biggl |\sum_{s^{\prime}} T(s'|s, a)  \max_{a^{\prime}} Q_0(gs^{\prime}, ga^{\prime}) \\
& \quad - \sum_{s^{\prime}}  T(gs'|gs, ga) \max_{a'} Q_0(gs', ga')\Biggl | \\
\end{aligned}
\end{equation*}}

If $\sum_{s^{\prime}} T(s'|s, a)  \max_{a^{\prime}} Q_0(s^{\prime}, a^{\prime})$  is greater than or equal to $\sum_{s^{\prime}} T(gs'|gs, ga)  \max_{a^{\prime}} Q_0(gs^{\prime}, ga^{\prime})$, then we have:
{\scriptsize
\begin{equation*}
\begin{aligned}
&   \Biggl |\sum_{s^{\prime}} T(s'|s, a)  \max_{a^{\prime}} Q_0(s^{\prime}, a^{\prime})  -  \sum_{s^{\prime}} T(gs'|gs, ga)  \max_{a^{\prime}} Q_0(gs^{\prime}, ga^{\prime}) \Biggl | \\
& \leq \Biggl |\sum_{s^{\prime}} T(s'|s, a)  \max_{a^{\prime}} (Q_0(gs^{\prime}, ga^{\prime})+\epsilon) \\
& \quad - \sum_{s^{\prime}}  T(gs'|gs, ga) \max_{a'} Q_0(gs', ga')\Biggl | \\
& \leq \epsilon  + \Biggl |\sum_{s^{\prime}} T(s'|s, a)  \max_{a^{\prime}} Q_0(gs^{\prime}, ga^{\prime}) \\
& \quad - \sum_{s^{\prime}}  T(gs'|gs, ga) \max_{a'} Q_0(gs', ga')\Biggl | \\
& = \epsilon + \Biggl | \mathbb{E}_{T(s'|s, a)}[f(s')] - \mathbb{E}_{T(gs'|gs, ga)}[f(s')] \Biggl |_{f(s')=\max_{a'} Q_0(gs', ga')}\\
& \leq \epsilon + \text{MMD}_{\mathcal{F}} \left(T(s'|s, a), T(gs'|gs, ga)\right) \\
& \leq \epsilon + \delta
\end{aligned}
\end{equation*}}

If $\sum_{s^{\prime}} T(s'|s, a)  \max_{a^{\prime}} Q_0(s^{\prime}, a^{\prime})$  is smaller than $\sum_{s^{\prime}} T(gs'|gs, ga)  \max_{a^{\prime}} Q_0(gs^{\prime}, ga^{\prime})$, then we have:

{\scriptsize
\begin{equation*}
\begin{aligned}
&   \Biggl |\sum_{s^{\prime}} T(s'|s, a)  \max_{a^{\prime}} Q_0(s^{\prime}, a^{\prime})  -  \sum_{s^{\prime}} T(gs'|gs, ga)  \max_{a^{\prime}} Q_0(gs^{\prime}, ga^{\prime}) \Biggl | \\
& = \Biggl | \sum_{s^{\prime}} T(gs'|gs, ga)  \max_{a^{\prime}} Q_0(gs^{\prime}, ga^{\prime})  -  \sum_{s^{\prime}} T(s'|s, a)  \max_{a^{\prime}} Q_0(s^{\prime}, a^{\prime}) \Biggl | \\
& \leq \Biggl | \sum_{s^{\prime}} T(gs'|gs, ga)  \max_{a^{\prime}} (Q_0(s^{\prime}, a^{\prime}) +\epsilon) \\
& \quad  -  \sum_{s^{\prime}} T(s'|s, a)  \max_{a^{\prime}} Q_0(s^{\prime}, a^{\prime}) \Biggl | \\
& \leq \epsilon  + \Biggl | \sum_{s^{\prime}} T(gs'|gs, ga)  \max_{a^{\prime}} Q_0(s^{\prime}, a^{\prime}) \\
& \quad  -  \sum_{s^{\prime}} T(s'|s, a)  \max_{a^{\prime}} Q_0(s^{\prime}, a^{\prime}) \Biggl | \\
& = \epsilon + \Biggl | \mathbb{E}_{T(gs'|gs, ga)}[f(s')] - \mathbb{E}_{T(s'|s, a)}[f(s')]  \Biggl |_{f(s')=\max_{a'} Q_0(s', a')}\\
& \leq \epsilon + \text{MMD}_{\mathcal{F}} \left(T(s'|s, a), T(gs'|gs, ga)\right) \\
& \leq \epsilon + \delta
\end{aligned}
\end{equation*}}

Taking the above discussions into account, we conclude that
{\scriptsize
\begin{equation*}
\begin{aligned}
& \left|Q_1(s, a) - Q_1(gs, ga)\right| \leq \epsilon +\gamma (\epsilon + \delta) = (1+\gamma)\epsilon+\gamma \delta
\end{aligned}
\end{equation*}}
\noindent
\textbf{Step 3 (induction step)}: As we have the base cases, let's assume that
{\scriptsize
\begin{equation*}
|Q_j(s, a) - Q_j(gs, ga)| \leq \sum_{i=0}^{j} \gamma^i\epsilon +\sum_{i=1}^{j}\gamma^i \delta
\end{equation*}}
holds for all non-negative integer $j<m$ and all state-action pairs $(s,a) \in\Psi$. Now we prove by induction on $m$ that the proposition is true. For the case $m$, we have: 
{\scriptsize
\begin{equation*}
\begin{aligned}
& \left|Q_m(s, a) - Q_m(gs, ga)\right| \\
& = \left| R(s, a) + \gamma \sum_{s'} T(s'|s, a)  \max_{a^{\prime}} Q_{m-1}(s', a') \right.\\
& \quad \left. - R(gs, ga) - \gamma \sum_{s'} T(gs'|gs, ga)  \max_{a^{\prime}} Q_{m-1}(gs^{\prime}, ga^{\prime}) \right| \\
& \leq \left| R(s, a) -  R(gs, ga) \right| +\gamma \biggl |\sum_{s'} T(s'|s, a) \max_{a^{\prime}} Q_{m-1}(s', a')  \\
& \quad - \sum_{s'} T(gs^{\prime}|gs, ga) \max_{a^{\prime}} Q_{m-1}(gs^{\prime}, ga^{\prime}) \biggl | \\
& \leq \epsilon + \gamma \biggl |\sum_{s'} T(s'|s, a) \max_{a^{\prime}} Q_{m-1}(s', a')  \\
& \quad - \sum_{s'} T(gs^{\prime}|gs, ga) \max_{a^{\prime}} Q_{m-1}(gs^{\prime}, ga^{\prime}) \biggl | \\
\end{aligned}
\end{equation*}}
By repeating the same discussions in Step 2 regarding two situations of 
$\sum_{s^{\prime}} T(s'|s, a)  \max_{a^{\prime}} Q_{m-1}(s^{\prime}, a^{\prime})$  and $\sum_{s^{\prime}} T(gs'|gs, ga)  \max_{a^{\prime}} Q_{m-1}(gs^{\prime}, ga^{\prime})$, we can obtain that:


{\scriptsize
\begin{equation*}
\begin{aligned}
&  \biggl |\sum_{s'} T(s'|s, a) \max_{a^{\prime}} Q_{m-1}(s', a')  \\
& \quad - \sum_{s'} T(gs^{\prime}|gs, ga) \max_{a^{\prime}} Q_{m-1}(gs^{\prime}, ga^{\prime}) \biggl | \\
& \leq \left( \sum_{i=0}^{m-1} \gamma^i\epsilon + \sum_{i=1}^{m-1}\gamma^i \delta \right) + \delta
\end{aligned}
\end{equation*}}

Therefore, we have

{\scriptsize
\begin{equation*}
\begin{aligned}
& \left|Q_m(s, a) - Q_m(gs, ga)\right| \\
& \leq \epsilon +\gamma \left(\sum_{i=0}^{m-1} \gamma^i\epsilon + \sum_{i=1}^{m-1}\gamma^i \delta\right) +\gamma \delta\\
& = \sum_{i=0}^{m} \gamma^i\epsilon + \sum_{i=1}^{m}\gamma^i \delta
\end{aligned}
\end{equation*}}

Based on mathematical induction, since the statement of the proposition in both the base step and  induction step has been proved as true, $|Q_{m}(s, a) - Q_{m}(gs, ga)|$ is bounded for all $m$. When $0<\gamma <1$ and $m$ approaches infinity, we have:

{\scriptsize
\begin{equation*}
\begin{aligned}
& \textit{Error}_{\mathcal{M}_g}=|Q^{\star}(s, a) - Q^{\star}(gs, ga)| \\
& = \lim_{{m \to \infty}} \Biggl |Q^{m}(s, a) - Q^{m}(gs, ga)\Biggl | \\
    &\leq\lim_{{m \to \infty}} \left(\sum_{{i=0}}^{m} \gamma^i\epsilon + \sum_{{i=1}}^{m}\gamma^i \delta\right) \\
    & = \frac{\epsilon}{1-\gamma}+\frac{\gamma \delta}{1-\gamma}
\end{aligned}
\end{equation*}}

The proof of proposition 1 is completed.
\end{proof}






\end{document}
